% GENERATED FILE. Edit canonical lessons, then run npm run generate:books.
% canonical-source-hash: fnv1a64:336fa41a7af32cc1
% canonical-lessons: ML-C493-upabhoktavu, ML-C493-itapatukaran, ML-C493-varanri, ML-C493-pariksikkuka, ML-C493-marriyetukkuka

\chapter{Buyers and Warranties}
\label{ch:ml-buyers-and-warranties}
\hlchaptermodality{493}

\begin{chapteropening}
\textbf{By the end of this chapter:} \emph{I can say a consumer, a customer, a warranty, to test and to swap for another.}

Complete the last lesson of chapter 493: I can say a consumer, a customer, a warranty, to test and to swap for another.
\end{chapteropening}

\section[\texorpdfstring{\ml{ഉപഭോക്താവ്} (upabhōktāvŭ)}{ഉപഭോക്താവ് (upabhōktāvŭ)}]{\ml{ഉപഭോക്താവ്} (upabhōktāvŭ) — a consumer}

\label{lesson:ML-C493-upabhoktavu}



\begin{quote}
Before the new one: say the Malayalam for genuine, then the Malayalam for to adulterate.
\end{quote}

\subsection*{You'll want to know: \ml{ഉപഭോക്താവ്}}
\textbf{\ml{ഉപഭോക്താവ്}} — \emph{upabhōktāvŭ} — \textquotedblleft{}a consumer\textquotedblright{}.

\begin{quote}
\ml{ഉപഭോക്താവിന്} \ml{ബിൽ} \ml{ചോദിക്കാൻ} \ml{അവകാശമുണ്ട്}.
\end{quote}

\begin{quote}
\emph{upabhōktāvinŭ bil cōdikkān avakāśamuṇṭŭ.}
\end{quote}

\begin{quote}
A consumer has the right to ask for a bill.
\end{quote}

A formal word; a shop calls its buyer a \ml{കസ്റ്റമർ}.

\begin{cousinweb}
\textbf{\ml{ഉപഭോക്താവ്}} (\emph{upabhōktāvŭ}) is borrowed from Sanskrit \emph{upabhoktṛ}, \textquotedblleft{}one who enjoys or makes use of\textquotedblright{} (Monier-Williams).
\end{cousinweb}

\begin{grammarlens}[title={What you've built}]
One more word for getting around.
\end{grammarlens}

\subsection*{Your turn}
\begin{itemize}
\raggedright
  \item \emph{Say it:} \emph{upabhōktāvŭ}
  \item \emph{Say it:} \emph{upabhōktāvŭ}, once more
  \item \emph{Say it:} \emph{upabhōktāvinŭ bil cōdikkān avakāśamuṇṭŭ.} — A consumer has the right to ask for a bill.
  \item \emph{Recall:} say the Malayalam for genuine, then the Malayalam for to adulterate, then say \emph{upabhōktāvŭ} again
\end{itemize}

\begin{tcolorbox}[breakable,colback=teal!4,colframe=teal!35!black,title={Before you move on}]
What does \ml{ഉപഭോക്താവ്} mean? (\textquotedblleft{}A consumer\textquotedblright{}.) Say it once more.
\end{tcolorbox}
\section[\texorpdfstring{\ml{ഇടപാടുകാരൻ} (iṭapāṭukāran)}{ഇടപാടുകാരൻ (iṭapāṭukāran)}]{\ml{ഇടപാടുകാരൻ} (iṭapāṭukāran) — a customer; a client}

\label{lesson:ML-C493-itapatukaran}



\begin{quote}
Before the new one: say the Malayalam for to adulterate, then the Malayalam for a consumer.
\end{quote}

\subsection*{You'll want to know: \ml{ഇടപാടുകാരൻ}}
\textbf{\ml{ഇടപാടുകാരൻ}} — \emph{iṭapāṭukāran} — \textquotedblleft{}a customer; a client\textquotedblright{}.

\begin{quote}
\ml{ആ} \ml{കടയിൽ} \ml{എപ്പോഴും} \ml{ധാരാളം} \ml{ഇടപാടുകാർ} \ml{ഉണ്ട്}.
\end{quote}

\begin{quote}
\emph{ā kaṭayil eppōḻuṁ dhārāḷaṁ iṭapāṭukār uṇṭŭ.}
\end{quote}

\begin{quote}
That shop has plenty of customers all the time.
\end{quote}

Plural: \ml{ഇടപാടുകാർ}, \emph{iṭapāṭukār}.

\ml{ഇടപാട്} is a dealing or a transaction; -\ml{കാരൻ} is the one who has it. Plural: \ml{ഇടപാടുകാർ}.

\begin{cousinweb}
\textbf{\ml{ഇടപാടുകാരൻ}} (\emph{iṭapāṭukāran}) is built from \ml{ഇടപാട്}, a dealing or transaction, with -\ml{കാരൻ}, the one who does it; the origin of \ml{ഇടപാട്} is not settled.
\end{cousinweb}

\begin{grammarlens}[title={What you've built}]
One more word for getting around.
\end{grammarlens}

\subsection*{Your turn}
\begin{itemize}
\raggedright
  \item \emph{Say it:} \emph{iṭapāṭukāran}
  \item \emph{Say it:} \emph{iṭapāṭukāran}, once more
  \item \emph{Say it:} \emph{ā kaṭayil eppōḻuṁ dhārāḷaṁ iṭapāṭukār uṇṭŭ.} — That shop has plenty of customers all the time.
  \item \emph{Recall:} say the Malayalam for to adulterate, then the Malayalam for a consumer, then say \emph{iṭapāṭukāran} again
\end{itemize}

\begin{tcolorbox}[breakable,colback=teal!4,colframe=teal!35!black,title={Before you move on}]
What does \ml{ഇടപാടുകാരൻ} mean? (\textquotedblleft{}A customer; a client\textquotedblright{}.) Say it once more.
\end{tcolorbox}
\section[\texorpdfstring{\ml{വാറന്റി} (vāṟanṟi)}{വാറന്റി (vāṟanṟi)}]{\ml{വാറന്റി} (vāṟanṟi) — a warranty}

\label{lesson:ML-C493-varanri}



\begin{quote}
Before the new one: say the Malayalam for a consumer, then the Malayalam for a customer.
\end{quote}

\subsection*{You'll want to know: \ml{വാറന്റി}}
\textbf{\ml{വാറന്റി}} — \emph{vāṟanṟi} — \textquotedblleft{}a warranty\textquotedblright{}.

\begin{quote}
\ml{ആ} \ml{ഫോണിന്} \ml{ഒരു} \ml{വർഷത്തെ} \ml{വാറന്റി} \ml{ഉണ്ട്}.
\end{quote}

\begin{quote}
\emph{ā phōṇinŭ oru varṣatte vāṟanṟi uṇṭŭ.}
\end{quote}

\begin{quote}
That phone has a one-year warranty.
\end{quote}

Keep the bill: the shop asks for it when you claim the warranty.

\begin{cousinweb}
\textbf{\ml{വാറന്റി}} (\emph{vāṟanṟi}) was borrowed from English \emph{warranty}, a maker's promise to repair a fault.
\end{cousinweb}

\begin{grammarlens}[title={What you've built}]
One more word for getting around.
\end{grammarlens}

\subsection*{Your turn}
\begin{itemize}
\raggedright
  \item \emph{Say it:} \emph{vāṟanṟi}
  \item \emph{Say it:} \emph{vāṟanṟi}, once more
  \item \emph{Say it:} \emph{ā phōṇinŭ oru varṣatte vāṟanṟi uṇṭŭ.} — That phone has a one-year warranty.
  \item \emph{Recall:} say the Malayalam for a consumer, then the Malayalam for a customer, then say \emph{vāṟanṟi} again
\end{itemize}

\begin{tcolorbox}[breakable,colback=teal!4,colframe=teal!35!black,title={Before you move on}]
What does \ml{വാറന്റി} mean? (\textquotedblleft{}A warranty\textquotedblright{}.) Say it once more.
\end{tcolorbox}
\section[\texorpdfstring{\ml{പരീക്ഷിക്കുക} (parīkṣikkuka)}{പരീക്ഷിക്കുക (parīkṣikkuka)}]{\ml{പരീക്ഷിക്കുക} (parīkṣikkuka) — to test; to try out}

\label{lesson:ML-C493-pariksikkuka}



\begin{quote}
Before the new one: say the Malayalam for a customer, then the Malayalam for a warranty.
\end{quote}

\subsection*{You'll want to know: \ml{പരീക്ഷിക്കുക}}
\textbf{\ml{പരീക്ഷിക്കുക}} — \emph{parīkṣikkuka} — \textquotedblleft{}to test; to try out\textquotedblright{}.

\begin{quote}
\ml{വാങ്ങുന്നതിന്} \ml{മുൻപ്} \ml{അവൻ} \ml{ഫാൻ} \ml{പരീക്ഷിച്ചു}.
\end{quote}

\begin{quote}
\emph{vāṅṅunnatinŭ munpŭ avan phān parīkṣiccu.}
\end{quote}

\begin{quote}
Before buying, he tested the fan.
\end{quote}

Past: \ml{പരീക്ഷിച്ചു}, \emph{parīkṣiccu}.

The verb of \ml{പരീക്ഷ}, an exam: to put something to the test.

\begin{cousinweb}
\textbf{\ml{പരീക്ഷിക്കുക}} (\emph{parīkṣikkuka}) is built from Sanskrit \emph{parīkṣā}, \textquotedblleft{}inspection, test\textquotedblright{} (CDIAL 7904), with the verb ending -\ml{ിക്കുക}.
\end{cousinweb}

\begin{grammarlens}[title={What you've built}]
One more word for getting around.
\end{grammarlens}

\subsection*{Your turn}
\begin{itemize}
\raggedright
  \item \emph{Say it:} \emph{parīkṣikkuka}
  \item \emph{Say it:} \emph{parīkṣikkuka}, once more
  \item \emph{Say it:} \emph{vāṅṅunnatinŭ munpŭ avan phān parīkṣiccu.} — Before buying, he tested the fan.
  \item \emph{Recall:} say the Malayalam for a customer, then the Malayalam for a warranty, then say \emph{parīkṣikkuka} again
\end{itemize}

\begin{tcolorbox}[breakable,colback=teal!4,colframe=teal!35!black,title={Before you move on}]
What does \ml{പരീക്ഷിക്കുക} mean? (\textquotedblleft{}To test; to try out\textquotedblright{}.) Say it once more.
\end{tcolorbox}
\section[\texorpdfstring{\ml{മാറ്റിയെടുക്കുക} (māṟṟiyeṭukkuka)}{മാറ്റിയെടുക്കുക (māṟṟiyeṭukkuka)}]{\ml{മാറ്റിയെടുക്കുക} (māṟṟiyeṭukkuka) — to swap for another}

\label{lesson:ML-C493-marriyetukkuka}



\begin{quote}
Before the new one: say the Malayalam for a warranty, then the Malayalam for to test.
\end{quote}

\subsection*{You'll want to know: \ml{മാറ്റിയെടുക്കുക}}
\textbf{\ml{മാറ്റിയെടുക്കുക}} — \emph{māṟṟiyeṭukkuka} — \textquotedblleft{}to swap for another\textquotedblright{}.

\begin{quote}
\ml{ചെറിയ} \ml{ഷർട്ട്} \ml{ഞാൻ} \ml{കടയിൽ} \ml{പോയി} \ml{മാറ്റിയെടുത്തു}.
\end{quote}

\begin{quote}
\emph{ceṟiya ṣarṭṭŭ ñān kaṭayil pōyi māṟṟiyeṭuttu.}
\end{quote}

\begin{quote}
I went to the shop and swapped the small shirt.
\end{quote}

Past: \ml{മാറ്റിയെടുത്തു}, \emph{māṟṟiyeṭuttu}.

\ml{മാറ്റി}, having changed, plus \ml{എടുക്കുക}, to take: to change a thing and take the new one.

\begin{cousinweb}
\textbf{\ml{മാറ്റിയെടുക്കുക}} (\emph{māṟṟiyeṭukkuka}) is a native Dravidian compound of \ml{മാറ്റുക}, to change (DEDR 4834), and \ml{എടുക്കുക}, to take (DEDR 851).
\end{cousinweb}

\begin{grammarlens}[title={What you've built}]
One more word for getting around.
\end{grammarlens}

\subsection*{Your turn}
\begin{itemize}
\raggedright
  \item \emph{Say it:} \emph{māṟṟiyeṭukkuka}
  \item \emph{Say it:} \emph{māṟṟiyeṭukkuka}, once more
  \item \emph{Say it:} \emph{ceṟiya ṣarṭṭŭ ñān kaṭayil pōyi māṟṟiyeṭuttu.} — I went to the shop and swapped the small shirt.
  \item \emph{Recall:} say the Malayalam for a warranty, then the Malayalam for to test, then say \emph{māṟṟiyeṭukkuka} again
\end{itemize}

\begin{tcolorbox}[breakable,colback=teal!4,colframe=teal!35!black,title={Before you move on}]
What does \ml{മാറ്റിയെടുക്കുക} mean? (\textquotedblleft{}To swap for another\textquotedblright{}.) Say it once more.
\end{tcolorbox}
